\section{Overview}
\label{sec:overview}

\subsection{Motivating Example}
\label{sec:motivating_example}

A network ingress path buffers packets in a ring so that bursts do not cause
drops. Bursts reach 800 packets, and the developer provisions 25\% headroom, so
the ring needs at least 1000 slots. The queue keeps two free-running 64-bit
counters. \texttt{tail} counts enqueued packets and \texttt{head} counts
dequeued ones. The queue holds \texttt{tail - head} packets. Each counter value
maps to a storage slot by taking it modulo the capacity.

\begin{lstlisting}[language=C]
#define CAPACITY 1000
uint64_t ring[CAPACITY];
uint64_t head = 0, tail = 0;

bool enqueue(const uint64_t pkt) {
    if (tail - head == CAPACITY) return false;
    ring[tail % CAPACITY] = pkt;
    ++tail;
    return true;
}
\end{lstlisting}

Every enqueue and dequeue computes a remainder by 1000. MSVC~19.29 at
\texttt{/O2} on x86-64 compiles that remainder to a reciprocal multiply:

\begin{lstlisting}[language={[x86masm]Assembler}]
; i % 1000
mov  rax, 442721857769029239
mul  rcx
mov  rax, rcx
sub  rax, rdx
shr  rax, 1
add  rax, rdx
shr  rax, 9
imul rax, rax, 1000
sub  rcx, rax
mov  rax, rcx
ret
\end{lstlisting}

A 1024-cell ring uses \texttt{i \% 1024}, which is the same as
\texttt{i \& 1023}, and the same compiler emits a mask:

\begin{lstlisting}[language={[x86masm]Assembler}]
; i % 1024
and  ecx, 1023
mov  eax, ecx
ret
\end{lstlisting}

Table~\ref{tab:motivating} splits the requirement, the source, what an as-if
compiler may emit, and what a contract admits. As-if will not rewrite 1000 to
1024, because the two queues disagree on the 1001st consecutive enqueue from
empty. In isolation, on an Intel Core i7-14700K, a dependent chain of slot
computations takes a median of 2.67\,ns per step with capacity 1000 and
0.39\,ns with capacity 1024, a factor of 6.8. When the computations are
independent, the factor is 3.1. These figures measure the slot computation
alone. Whole queues, where memory traffic dilutes the difference, are measured
in Section~\ref{sec:evaluation}, and the larger ring uses 24 more slots.

Both capacities meet a requirement of at least 1000 slots. Which is better
depends on the deployment. On a server where the path is latency-critical and
memory is plentiful, 1024 is the better choice. On a microcontroller whose
memory budget is exactly the 1000-slot footprint, the extra 24 slots do not
fit. The preference is specific to the target, so neither capacity is right
in the source as a universal constant.

\begin{table}
\caption{Columns are the requirement, the source \(P_1\), what MSVC
\texttt{/O2} emits under as-if, and a program the contract admits. As-if
must keep the remainder that is in the source. A contract with
\texttt{CAPACITY}~\(\ge 1000\) admits 1024.}
\label{tab:motivating}
\scriptsize
\begin{tabular}{p{0.22\linewidth}p{0.22\linewidth}p{0.22\linewidth}p{0.22\linewidth}}
\toprule
Requirement & Source \(P_1\) & MSVC \texttt{/O2} (as-if) & Contract admits \\
\midrule
\(\texttt{CAPACITY} \ge 1000\), FIFO, full at \texttt{CAPACITY} &
\texttt{\#define CAPACITY 1000}; \texttt{i \% CAPACITY} &
reciprocal multiply for \texttt{\% 1000}; 1001st enqueue fails &
\texttt{CAPACITY 1024}; \texttt{i \& 1023}; 1001st enqueue succeeds \\
full at 1000; backing store not in the API &
\texttt{CAPACITY} and \texttt{STORAGE} both 1000; \texttt{i \% STORAGE} &
keeps remainder-by-1000 &
\texttt{STORAGE 1024}, \texttt{CAPACITY 1000}; mask; enqueue still fails at 1000 \\
\bottomrule
\end{tabular}
\end{table}

That change is not the only way to use a power-of-two backing store. A library
can allocate 1024 slots and still report full at 1000, using the extra slots
only so that indexing is a mask. Folly's queue after September 2026 does
exactly that: it rounds storage up internally and keeps the caller-facing
capacity~\cite{folly_pcq}. The two designs use the same hardware trick and are different contracts: one
changes when enqueue fails, and the other does not.

\subsection{What Existing Tools Can Do}
\label{sec:existing}

\paragraph{Conventional Compilers and Translation Validation.}
The two queues that become full at different lengths are not observationally
equivalent. Translation validators such as Alive2~\cite{lopes2021alive2} check
refinement against the source, which allows extra freedom only where the source
has undefined behavior, and this queue has none on its intended inputs.
Alive2 rejects the 1024-cell ring against the 1000-cell source: the two
disagree on the 1001st enqueue from empty.

\paragraph{Compiler Flags.}
Production compilers already relax as-if under flags, and
\texttt{-ffast-math} is one example: it licenses reassociation and other
algebraic rewrites that IEEE-754 would forbid, and the waiver can be scoped
with pragmas. Those rewrites change observable numeric behavior, including
treating NaN and infinities as unused, sometimes with measurable
implications, but they do not let the compiler replace a 1000-cell remainder
with a 1024-cell mask.

\paragraph{Compiler Data-Layout Transformations.}
Production compilers do change data layout. IBM's XL compilers reshape arrays
of structures~\cite{zhao2007forma}, HP's compilers split and reorder
structures~\cite{hundt2006practical}, and automatic pool allocation regroups
heap objects~\cite{lattner2005pool}. These transformations must remain
invisible to the program. They are often abandoned when the analysis cannot
prove that~\cite{hundt2006practical}, and changing a logical capacity is
visible by construction, so it is outside their reach. Padding a private array while
preserving the API is closer to what they already do, and still requires
proving that the extra cells are never observed.

\paragraph{Annotations and Assumptions.}
C++'s \texttt{[[assume]]} states facts the compiler may rely on. A false
assumption is undefined behavior, and the C++26 contracts design leaves out an
assume semantic for that reason~\cite{p2899r1}. Recent Clang attributes permit
array-of-structures to structure-of-arrays conversions on annotated
code~\cite{radtke2025annotation}. Both are trusted rather than checked, and
neither states that a capacity may grow.

\paragraph{Scheduling Languages.}
Halide and Exo separate an algorithm from its schedule~\cite{ragan2013halide,
ikarashi2022exo}. The algorithm's meaning, including the size of the ring,
stays fixed.

\paragraph{Specification-Based Synthesis.}
Cozy synthesizes collection implementations from specifications, and Primrose
selects container implementations that satisfy stated
properties~\cite{loncaric2016cozy,Qin2023Primrose}. Fiat refines specifications
into verified implementations in Coq~\cite{delaware2015fiat}. These systems
accept freedom from a specification, in their own languages or libraries.
Primrose comes closest to existing code: it reads annotated Rust and selects a
library implementation. This paper starts from an existing C program, a
contract over its observable behavior, and that language's production compiler.

\subsection{Intent Compilation on the Example}
\label{sec:overview_intent}

The developer states the requirement instead of a single capacity. In
\textsc{IntC} that is a YAML file plus C harnesses. The contract has four pieces:
a constraint on named parameters, a functional harness, an optional resource
budget, and a measured objective. On this queue:

\begin{lstlisting}[language=C]
/* CAPACITY >= 1000 */
/* FIFO; enqueue fails iff the ring holds CAPACITY packets */
/* optional static-memory ceiling */
/* measured ns/op on a given target */
\end{lstlisting}

\textsc{IntentCompile} (\S\ref{sec:pipeline}) rewrites \texttt{e \% CAPACITY} to
\texttt{((e) \& (CAPACITY - 1))}, which requires a power-of-two capacity, searches
for the admissible power of two closest to 1000, and asks a bounded model
checker whether every sequence of at most \(K\) operations from an empty queue
satisfies the harness (the bound of \(\models_K\), \S\ref{sec:verifier}).
Under a server profile with no tight memory budget, 1024
is admitted and, if it measures faster, emitted as source. Under an 8016-byte
static-memory ceiling (the 1000-slot baseline's footprint), 1024 fails the
budget at 8208 bytes and the baseline is kept.
It also tries a pad-and-mask form: a \texttt{STORAGE} hole, a mask on that
store, and fullness still at \texttt{CAPACITY}. That candidate is the
unexposed-slack case of Section~\ref{sec:slack}. When the baseline already names
both parameters, the mask rewrite on \texttt{STORAGE} is enough.
